\documentclass[12pt,a4paper]{article}
\usepackage[utf8]{inputenc}
\usepackage[T5]{fontenc}          % Bắt buộc T5 cho tiếng Việt trong pdflatex
\usepackage[vietnamese]{babel}
\usepackage{amsmath, amssymb, amsfonts, amsthm}
\usepackage{geometry}
\usepackage{enumitem}

\geometry{margin=2cm}

\newcommand{\R}{\mathbb{R}}
\newcommand{\N}{\mathbb{N}}
\newcommand{\ip}[2]{\left\langle #1,#2\right\rangle}
\newcommand{\norm}[1]{\left\|#1\right\|}

\begin{document}



\subsection*{BÀI 1A (Mẫu Câu 1 trên $V = \R^5$):}
Trên $V=\R^5$, cho tập hợp
\[
W = \left\{ X = (x_1; x_2; x_3; x_4; x_5) \in \R^5 \left| \begin{matrix} x_1 - x_2 + 2x_3 - x_4 + x_5 = 0 \\ 2x_1 - x_2 + 3x_3 - 2x_4 + x_5 = 0 \end{matrix} \right. \right\}
\]
Chứng minh rằng $W$ là không gian véc tơ con của $V$. Tìm hệ sinh, cơ sở và xác định số chiều cho $W$.

\underline{\textbf{Giải:}}

Từ điều kiện của $W$, ta viết lại hệ phương trình tuyến tính thuần nhất dưới dạng ma trận hóa như sau:
\[
\begin{pmatrix} 
1 & -1 & 2 & -1 & 1 \mid 0 \\ 
2 & -1 & 3 & -2 & 1 \mid 0 
\end{pmatrix}
\xrightarrow{(2) \to (2) - 2(1)}
\begin{pmatrix} 
1 & -1 & 2 & -1 & 1 \mid 0 \\ 
0 & 1 & -1 & 0 & -1 \mid 0 
\end{pmatrix}
\]
Do cột 3, 4, 5 không bán chuẩn hóa được nên hệ có vô số nghiệm như sau:

Đặt $\begin{cases} x_3 = a \\ x_4 = b \\ x_5 = c \end{cases} \quad \forall a, b, c \in \R$.

Từ dòng (2) ta có: $x_2 - x_3 - x_5 = 0 \Rightarrow x_2 = x_3 + x_5 = a + c$.

Từ dòng (1) ta có: $x_1 - x_2 + 2x_3 - x_4 + x_5 = 0 \Rightarrow x_1 = x_2 - 2x_3 + x_4 - x_5 = (a+c) - 2a + b - c = -a + b$.

Thay vào $W$ ta được:
\begin{align*}
W &= \{ X = (-a+b; a+c; a; b; c) \mid a, b, c \in \R \} \\
  &= \{ X = (-a; a; a; 0; 0) + (b; 0; 0; b; 0) + (0; c; 0; 0; c) \mid a, b, c \in \R \} \\
  &= \{ X = a(-1; 1; 1; 0; 0) + b(1; 0; 0; 1; 0) + c(0; 1; 0; 0; 1) \mid a, b, c \in \R \}.
\end{align*}
\begin{quote}
\small\textbf{(* Note giải thích bước tách tọa độ bằng bảng hệ số theo $a, b, c$):}
\begin{center}
\renewcommand{\arraystretch}{1.15}
\setlength{\tabcolsep}{6pt}
\footnotesize
\begin{tabular}{|c|c|c|c|c|}
\hline
\textbf{Vị trí} & \textbf{Biểu thức tọa độ} & \textbf{Hệ số $a$} & \textbf{Hệ số $b$} & \textbf{Hệ số $c$} \\
\hline
Tọa độ 1 & $-a + b$ & $\mathbf{-1}$ & $\mathbf{1}$ & $\mathbf{0}$ \\
\hline
Tọa độ 2 & $a + c$ & $\mathbf{1}$ & $\mathbf{0}$ & $\mathbf{1}$ \\
\hline
Tọa độ 3 & $a$ & $\mathbf{1}$ & $\mathbf{0}$ & $\mathbf{0}$ \\
\hline
Tọa độ 4 & $b$ & $\mathbf{0}$ & $\mathbf{1}$ & $\mathbf{0}$ \\
\hline
Tọa độ 5 & $c$ & $\mathbf{0}$ & $\mathbf{0}$ & $\mathbf{1}$ \\
\hline
\end{tabular}
\end{center}
\textbf{Đọc các hệ số theo từng cột dọc từ trên xuống dưới:}
\begin{itemize}[noitemsep,topsep=2pt]
\item \textbf{Cột $a$:} Ta được véc tơ $\alpha_1 = (-1; 1; 1; 0; 0) \implies$ phần chứa $a$ là $a\alpha_1$.
\item \textbf{Cột $b$:} Ta được véc tơ $\alpha_2 = (1; 0; 0; 1; 0) \implies$ phần chứa $b$ là $b\alpha_2$.
\item \textbf{Cột $c$:} Ta được véc tơ $\alpha_3 = (0; 1; 0; 0; 1) \implies$ phần chứa $c$ là $c\alpha_3$.
\end{itemize}
$\Rightarrow$ Tách thành $X = a\alpha_1 + b\alpha_2 + c\alpha_3$ cực kỳ nhanh và không bao giờ bị nhầm dấu!
\end{quote}

Đặt $\alpha_1 = (-1; 1; 1; 0; 0), \alpha_2 = (1; 0; 0; 1; 0), \alpha_3 = (0; 1; 0; 0; 1)$ và gọi $S = \{ \alpha_1, \alpha_2, \alpha_3 \}$.

$\Rightarrow W = \{ X = a\alpha_1 + b\alpha_2 + c\alpha_3 \mid a, b, c \in \R; \alpha_1, \alpha_2, \alpha_3 \in S \} = \langle S \rangle$.

Mà $\langle S \rangle \le V$ nên $W = \langle S \rangle \le V$ (đpcm).

Ta có $W = \langle S \rangle$ nên $S = \{ \alpha_1, \alpha_2, \alpha_3 \}$ là một hệ sinh của $W$.

Tiếp theo, ta tìm cơ sở và số chiều cho $W$ như sau:

Lập ma trận
\[
A_S = \begin{pmatrix} \alpha_1 \\ \alpha_2 \\ \alpha_3 \end{pmatrix} = \begin{pmatrix} -1 & 1 & 1 & 0 & 0 \\ 1 & 0 & 0 & 1 & 0 \\ 0 & 1 & 0 & 0 & 1 \end{pmatrix} 
\xrightarrow{(2) \to (2) + (1)} \begin{pmatrix} -1 & 1 & 1 & 0 & 0 \\ 0 & 1 & 1 & 1 & 0 \\ 0 & 1 & 0 & 0 & 1 \end{pmatrix} 
\]
\[
\xrightarrow{(3) \to (3) - (2)} \begin{pmatrix} -1 & 1 & 1 & 0 & 0 \\ 0 & 1 & 1 & 1 & 0 \\ 0 & 0 & -1 & -1 & 1 \end{pmatrix}
\]
Do ma trận đáp số có 3 dòng khác zero nên $r(A_S) = 3 = $ số véc tơ trong $S$, nên $S$ vừa là hệ sinh, vừa ĐLTT nên $S$ cũng là 1 cơ sở của $W$, và số chiều là $\dim_{\R} W = 3$.

\vspace{0.3cm}

\subsection*{BÀI 1B (Mẫu Câu 1 trên $V = \R^6$):}
Trên $V=\R^6$, cho tập hợp
\[
W = \left\{ X = (x_1; x_2; x_3; x_4; x_5; x_6) \in \R^6 \left| \begin{matrix} x_1 + 2x_2 - x_3 + x_4 - 2x_5 + x_6 = 0 \\ 2x_1 + 4x_2 - x_3 + 3x_4 - 3x_5 + 3x_6 = 0 \end{matrix} \right. \right\}
\]
Chứng minh rằng $W$ là không gian véc tơ con của $V$. Tìm hệ sinh, cơ sở và xác định số chiều cho $W$.

\underline{\textbf{Giải:}}

Từ điều kiện của $W$, ta viết lại hệ phương trình tuyến tính thuần nhất dưới dạng ma trận hóa như sau:
\[
\begin{pmatrix} 
1 & 2 & -1 & 1 & -2 & 1 \mid 0 \\ 
2 & 4 & -1 & 3 & -3 & 3 \mid 0 
\end{pmatrix}
\xrightarrow{(2) \to (2) - 2(1)}
\begin{pmatrix} 
1 & 2 & -1 & 1 & -2 & 1 \mid 0 \\ 
0 & 0 & 1 & 1 & 1 & 1 \mid 0 
\end{pmatrix}
\]
Do cột 2, 4, 5, 6 không bán chuẩn hóa được nên hệ có vô số nghiệm như sau:

Đặt $\begin{cases} x_2 = a \\ x_4 = b \\ x_5 = c \\ x_6 = d \end{cases} \quad \forall a, b, c, d \in \R$.

Từ dòng (2) ta có: $x_3 + x_4 + x_5 + x_6 = 0 \Rightarrow x_3 = -x_4 - x_5 - x_6 = -b - c - d$.

Từ dòng (1) ta có: $x_1 = -2x_2 + x_3 - x_4 + 2x_5 - x_6 = -2a + (-b - c - d) - b + 2c - d = -2a - 2b + c - 2d$.

\begin{align*}
W &= \{ X = (-2a - 2b + c - 2d; a; -b - c - d; b; c; d) \mid a, b, c, d \in \R \} \\
  &= \{ X = a(-2; 1; 0; 0; 0; 0) + b(-2; 0; -1; 1; 0; 0) \\
  &\qquad\quad + c(1; 0; -1; 0; 1; 0) + d(-2; 0; -1; 0; 0; 1) \mid a, b, c, d \in \R \}.
\end{align*}
\begin{quote}
\small\textbf{(* Note giải thích bước tách tọa độ bằng bảng hệ số theo $a, b, c, d$):}
\begin{center}
\renewcommand{\arraystretch}{1.15}
\setlength{\tabcolsep}{5pt}
\footnotesize
\begin{tabular}{|c|c|c|c|c|c|}
\hline
\textbf{Vị trí} & \textbf{Biểu thức tọa độ} & \textbf{Hệ số $a$} & \textbf{Hệ số $b$} & \textbf{Hệ số $c$} & \textbf{Hệ số $d$} \\
\hline
Tọa độ 1 & $-2a - 2b + 1c - 2d$ & $\mathbf{-2}$ & $\mathbf{-2}$ & $\mathbf{1}$ & $\mathbf{-2}$ \\
\hline
Tọa độ 2 & $1a$ & $\mathbf{1}$ & $\mathbf{0}$ & $\mathbf{0}$ & $\mathbf{0}$ \\
\hline
Tọa độ 3 & $-1b - 1c - 1d$ & $\mathbf{0}$ & $\mathbf{-1}$ & $\mathbf{-1}$ & $\mathbf{-1}$ \\
\hline
Tọa độ 4 & $1b$ & $\mathbf{0}$ & $\mathbf{1}$ & $\mathbf{0}$ & $\mathbf{0}$ \\
\hline
Tọa độ 5 & $1c$ & $\mathbf{0}$ & $\mathbf{0}$ & $\mathbf{1}$ & $\mathbf{0}$ \\
\hline
Tọa độ 6 & $1d$ & $\mathbf{0}$ & $\mathbf{0}$ & $\mathbf{0}$ & $\mathbf{1}$ \\
\hline
\end{tabular}
\end{center}
\textbf{Đọc các hệ số theo từng cột dọc từ trên xuống dưới:}
\begin{itemize}[noitemsep,topsep=2pt]
\item \textbf{Cột $a$:} Ta được véc tơ $\alpha_1 = (-2; 1; 0; 0; 0; 0) \implies$ phần chứa $a$ là $a\alpha_1$.
\item \textbf{Cột $b$:} Ta được véc tơ $\alpha_2 = (-2; 0; -1; 1; 0; 0) \implies$ phần chứa $b$ là $b\alpha_2$.
\item \textbf{Cột $c$:} Ta được véc tơ $\alpha_3 = (1; 0; -1; 0; 1; 0) \implies$ phần chứa $c$ là $c\alpha_3$.
\item \textbf{Cột $d$:} Ta được véc tơ $\alpha_4 = (-2; 0; -1; 0; 0; 1) \implies$ phần chứa $d$ là $d\alpha_4$.
\end{itemize}
$\Rightarrow$ Tách thành $X = a\alpha_1 + b\alpha_2 + c\alpha_3 + d\alpha_4$ cực kỳ nhanh và không bao giờ bị nhầm dấu!
\end{quote}
Đặt $\alpha_1 = (-2; 1; 0; 0; 0; 0)$, $\alpha_2 = (-2; 0; -1; 1; 0; 0)$, $\alpha_3 = (1; 0; -1; 0; 1; 0)$, $\alpha_4 = (-2; 0; -1; 0; 0; 1)$ và gọi $S = \{ \alpha_1, \alpha_2, \alpha_3, \alpha_4 \}$.

$\Rightarrow W = \{ X = a\alpha_1 + b\alpha_2 + c\alpha_3 + d\alpha_4 \mid a, b, c, d \in \R; \alpha_i \in S \} = \langle S \rangle$.

Mà $\langle S \rangle \le V$ nên $W = \langle S \rangle \le V$ (đpcm).

Ta có $W = \langle S \rangle$ nên $S = \{ \alpha_1, \alpha_2, \alpha_3, \alpha_4 \}$ là một hệ sinh của $W$.

Tiếp theo, ta tìm cơ sở và số chiều cho $W$ như sau:

Lập ma trận
\[
A_S = \begin{pmatrix} \alpha_1 \\ \alpha_2 \\ \alpha_3 \\ \alpha_4 \end{pmatrix} = 
\begin{pmatrix} 
-2 & 1 & 0 & 0 & 0 & 0 \\ 
-2 & 0 & -1 & 1 & 0 & 0 \\ 
1 & 0 & -1 & 0 & 1 & 0 \\ 
-2 & 0 & -1 & 0 & 0 & 1 
\end{pmatrix}
\xrightarrow{\substack{(2) \to (2) - (1) \\ (3) \to 2(3) + (1) \\ (4) \to (4) - (1)}}
\begin{pmatrix} 
-2 & 1 & 0 & 0 & 0 & 0 \\ 
0 & -1 & -1 & 1 & 0 & 0 \\ 
0 & 1 & -2 & 0 & 2 & 0 \\ 
0 & 0 & 0 & -1 & 0 & 1 
\end{pmatrix}
\]
\[
\xrightarrow{(3) \to (3) + (2)}
\begin{pmatrix} 
-2 & 1 & 0 & 0 & 0 & 0 \\ 
0 & -1 & -1 & 1 & 0 & 0 \\ 
0 & 0 & -3 & 1 & 2 & 0 \\ 
0 & 0 & 0 & -1 & 0 & 1 
\end{pmatrix}
\]
Do ma trận đáp số có 4 dòng khác zero nên $r(A_S) = 4 = $ số véc tơ trong $S$, nên $S$ vừa là hệ sinh, vừa ĐLTT nên $S$ cũng là 1 cơ sở của $W$, và số chiều là $\dim_{\R} W = 4$.

\vspace{0.3cm}

\subsection*{BÀI 1C (Mẫu Câu 1 trên $V = \R^6$ với hệ 3 phương trình ràng buộc):}
Trên $V=\R^6$, cho tập hợp
\[
W = \left\{ X = (x_1; x_2; x_3; x_4; x_5; x_6) \in \R^6 \left| \begin{matrix} x_1 - x_2 + x_3 - 2x_4 + x_5 - x_6 = 0 \\ 2x_1 - x_2 + 3x_3 - x_4 + 2x_5 - x_6 = 0 \\ x_1 + x_3 + 3x_4 + x_5 + x_6 = 0 \end{matrix} \right. \right\}
\]
Chứng minh rằng $W$ là không gian véc tơ con của $V$. Tìm hệ sinh, cơ sở và xác định số chiều cho $W$.

\underline{\textbf{Giải:}}

Từ điều kiện của $W$, ta viết lại hệ phương trình tuyến tính thuần nhất dưới dạng ma trận hóa như sau:
\[
\begin{pmatrix} 
1 & -1 & 1 & -2 & 1 & -1 \mid 0 \\ 
2 & -1 & 3 & -1 & 2 & -1 \mid 0 \\ 
1 & 0 & 1 & 3 & 1 & 1 \mid 0 
\end{pmatrix}
\xrightarrow{\substack{(2) \to (2) - 2(1) \\ (3) \to (3) - (1)}}
\begin{pmatrix} 
1 & -1 & 1 & -2 & 1 & -1 \mid 0 \\ 
0 & 1 & 1 & 3 & 0 & 1 \mid 0 \\ 
0 & 1 & 0 & 5 & 0 & 2 \mid 0 
\end{pmatrix}
\]
\[
\xrightarrow{(3) \to (3) - (2)}
\begin{pmatrix} 
1 & -1 & 1 & -2 & 1 & -1 \mid 0 \\ 
0 & 1 & 1 & 3 & 0 & 1 \mid 0 \\ 
0 & 0 & -1 & 2 & 0 & 1 \mid 0 
\end{pmatrix}
\xrightarrow{(3) \to -(3)}
\begin{pmatrix} 
1 & -1 & 1 & -2 & 1 & -1 \mid 0 \\ 
0 & 1 & 1 & 3 & 0 & 1 \mid 0 \\ 
0 & 0 & 1 & -2 & 0 & -1 \mid 0 
\end{pmatrix}
\]
Do cột 4, 5, 6 không bán chuẩn hóa được nên hệ có vô số nghiệm như sau:

Đặt $\begin{cases} x_4 = a \\ x_5 = b \\ x_6 = c \end{cases} \quad \forall a, b, c \in \R$.

Từ dòng (3) ta có: $x_3 - 2x_4 - x_6 = 0 \Rightarrow x_3 = 2x_4 + x_6 = 2a + c$.

Từ dòng (2) ta có: $x_2 + x_3 + 3x_4 + x_6 = 0 \Rightarrow x_2 = -x_3 - 3x_4 - x_6 = -(2a + c) - 3a - c = -5a - 2c$.

Từ dòng (1) ta có:
\begin{align*}
x_1 &= x_2 - x_3 + 2x_4 - x_5 + x_6 \\
    &= (-5a - 2c) - (2a + c) + 2a - b + c = -5a - b - 2c.
\end{align*}
Thay vào $W$ ta được:
\begin{align*}
W &= \{ X = (-5a - b - 2c; -5a - 2c; 2a + c; a; b; c) \mid a, b, c \in \R \} \\
  &= \{ X = a(-5; -5; 2; 1; 0; 0) + b(-1; 0; 0; 0; 1; 0) + c(-2; -2; 1; 0; 0; 1) \mid a, b, c \in \R \}.
\end{align*}
\begin{quote}
\small\textbf{(* Note giải thích bước tách tọa độ bằng bảng hệ số theo $a, b, c$):}
\begin{center}
\renewcommand{\arraystretch}{1.15}
\setlength{\tabcolsep}{6pt}
\footnotesize
\begin{tabular}{|c|c|c|c|c|}
\hline
\textbf{Vị trí} & \textbf{Biểu thức tọa độ} & \textbf{Hệ số $a$} & \textbf{Hệ số $b$} & \textbf{Hệ số $c$} \\
\hline
Tọa độ 1 & $-5a - 1b - 2c$ & $\mathbf{-5}$ & $\mathbf{-1}$ & $\mathbf{-2}$ \\
\hline
Tọa độ 2 & $-5a - 2c$ & $\mathbf{-5}$ & $\mathbf{0}$ & $\mathbf{-2}$ \\
\hline
Tọa độ 3 & $2a + 1c$ & $\mathbf{2}$ & $\mathbf{0}$ & $\mathbf{1}$ \\
\hline
Tọa độ 4 & $1a$ & $\mathbf{1}$ & $\mathbf{0}$ & $\mathbf{0}$ \\
\hline
Tọa độ 5 & $1b$ & $\mathbf{0}$ & $\mathbf{1}$ & $\mathbf{0}$ \\
\hline
Tọa độ 6 & $1c$ & $\mathbf{0}$ & $\mathbf{0}$ & $\mathbf{1}$ \\
\hline
\end{tabular}
\end{center}
\textbf{Đọc các hệ số theo từng cột dọc từ trên xuống dưới:}
\begin{itemize}[noitemsep,topsep=2pt]
\item \textbf{Cột $a$:} Ta được véc tơ $\alpha_1 = (-5; -5; 2; 1; 0; 0) \implies$ phần chứa $a$ là $a\alpha_1$.
\item \textbf{Cột $b$:} Ta được véc tơ $\alpha_2 = (-1; 0; 0; 0; 1; 0) \implies$ phần chứa $b$ là $b\alpha_2$.
\item \textbf{Cột $c$:} Ta được véc tơ $\alpha_3 = (-2; -2; 1; 0; 0; 1) \implies$ phần chứa $c$ là $c\alpha_3$.
\end{itemize}
$\Rightarrow$ Tách thành $X = a\alpha_1 + b\alpha_2 + c\alpha_3$ cực kỳ nhanh và không bao giờ bị nhầm dấu!
\end{quote}
Đặt $\alpha_1 = (-5; -5; 2; 1; 0; 0)$, $\alpha_2 = (-1; 0; 0; 0; 1; 0)$, $\alpha_3 = (-2; -2; 1; 0; 0; 1)$ và gọi $S = \{ \alpha_1, \alpha_2, \alpha_3 \}$.

$\Rightarrow W = \{ X = a\alpha_1 + b\alpha_2 + c\alpha_3 \mid a, b, c \in \R; \alpha_i \in S \} = \langle S \rangle$.

Mà $\langle S \rangle \le V$ nên $W = \langle S \rangle \le V$ (đpcm).

Ta có $W = \langle S \rangle$ nên $S = \{ \alpha_1, \alpha_2, \alpha_3 \}$ là một hệ sinh của $W$.

Tiếp theo, ta tìm cơ sở và số chiều cho $W$ như sau:

Lập ma trận
\[
A_S = \begin{pmatrix} \alpha_1 \\ \alpha_2 \\ \alpha_3 \end{pmatrix} = 
\begin{pmatrix} 
-5 & -5 & 2 & 1 & 0 & 0 \\ 
-1 & 0 & 0 & 0 & 1 & 0 \\ 
-2 & -2 & 1 & 0 & 0 & 1 
\end{pmatrix}
\xrightarrow{\substack{(1) \leftrightarrow (2) \\ (1) \to -(1)}}
\begin{pmatrix} 
1 & 0 & 0 & 0 & -1 & 0 \\ 
-5 & -5 & 2 & 1 & 0 & 0 \\ 
-2 & -2 & 1 & 0 & 0 & 1 
\end{pmatrix}
\]
\[
\xrightarrow{\substack{(2) \to (2) + 5(1) \\ (3) \to (3) + 2(1)}}
\begin{pmatrix} 
1 & 0 & 0 & 0 & -1 & 0 \\ 
0 & -5 & 2 & 1 & -5 & 0 \\ 
0 & -2 & 1 & 0 & -2 & 1 
\end{pmatrix}
\xrightarrow{(3) \to 5(3) - 2(2)}
\begin{pmatrix} 
1 & 0 & 0 & 0 & -1 & 0 \\ 
0 & -5 & 2 & 1 & -5 & 0 \\ 
0 & 0 & 1 & -2 & 0 & 5 
\end{pmatrix}
\]
Do ma trận đáp số có 3 dòng khác zero nên $r(A_S) = 3 = $ số véc tơ trong $S$, nên $S$ vừa là hệ sinh, vừa ĐLTT nên $S$ cũng là 1 cơ sở của $W$, và số chiều là $\dim_{\R} W = 3$.

\vspace{0.3cm}

\subsection*{BÀI 1D (Mẫu Câu 1 trên $V = \R^5$ với hệ chỉ có 1 phương trình ràng buộc):}
Trên $V=\R^5$, cho tập hợp
\[
W = \left\{ X = (x_1; x_2; x_3; x_4; x_5) \in \R^5 \mid x_1 - 2x_2 + 3x_3 - x_4 + 2x_5 = 0 \right\}
\]
Chứng minh rằng $W$ là không gian véc tơ con của $V$. Tìm hệ sinh, cơ sở và xác định số chiều cho $W$.

\underline{\textbf{Giải:}}

Từ điều kiện của $W$, ta viết lại phương trình tuyến tính thuần nhất dưới dạng ma trận hóa như sau:
\[
\begin{pmatrix} 1 & -2 & 3 & -1 & 2 \mid 0 \end{pmatrix}
\]
Ma trận đã ở dạng bậc thang. Do cột 2, 3, 4, 5 không bán chuẩn hóa được nên hệ có vô số nghiệm như sau:

Đặt $\begin{cases} x_2 = a \\ x_3 = b \\ x_4 = c \\ x_5 = d \end{cases} \quad \forall a, b, c, d \in \R$.

Từ phương trình ta có:
\[
x_1 = 2x_2 - 3x_3 + x_4 - 2x_5 = 2a - 3b + c - 2d.
\]
Thay vào $W$ ta được:
\begin{align*}
W &= \{ X = (2a - 3b + c - 2d; a; b; c; d) \mid a, b, c, d \in \R \} \\
  &= \{ X = a(2; 1; 0; 0; 0) + b(-3; 0; 1; 0; 0) + c(1; 0; 0; 1; 0) + d(-2; 0; 0; 0; 1) \mid a, b, c, d \in \R \}.
\end{align*}
\begin{quote}
\small\textbf{(* Note giải thích bước tách tọa độ bằng bảng hệ số theo $a, b, c, d$):}
\begin{center}
\renewcommand{\arraystretch}{1.15}
\setlength{\tabcolsep}{6pt}
\footnotesize
\begin{tabular}{|c|c|c|c|c|c|}
\hline
\textbf{Vị trí} & \textbf{Biểu thức tọa độ} & \textbf{Hệ số $a$} & \textbf{Hệ số $b$} & \textbf{Hệ số $c$} & \textbf{Hệ số $d$} \\
\hline
Tọa độ 1 & $2a - 3b + 1c - 2d$ & $\mathbf{2}$ & $\mathbf{-3}$ & $\mathbf{1}$ & $\mathbf{-2}$ \\
\hline
Tọa độ 2 & $1a$ & $\mathbf{1}$ & $\mathbf{0}$ & $\mathbf{0}$ & $\mathbf{0}$ \\
\hline
Tọa độ 3 & $1b$ & $\mathbf{0}$ & $\mathbf{1}$ & $\mathbf{0}$ & $\mathbf{0}$ \\
\hline
Tọa độ 4 & $1c$ & $\mathbf{0}$ & $\mathbf{0}$ & $\mathbf{1}$ & $\mathbf{0}$ \\
\hline
Tọa độ 5 & $1d$ & $\mathbf{0}$ & $\mathbf{0}$ & $\mathbf{0}$ & $\mathbf{1}$ \\
\hline
\end{tabular}
\end{center}
\textbf{Đọc các hệ số theo từng cột dọc từ trên xuống dưới:}
\begin{itemize}[noitemsep,topsep=2pt]
\item \textbf{Cột $a$:} Ta được véc tơ $\alpha_1 = (2; 1; 0; 0; 0) \implies$ phần chứa $a$ là $a\alpha_1$.
\item \textbf{Cột $b$:} Ta được véc tơ $\alpha_2 = (-3; 0; 1; 0; 0) \implies$ phần chứa $b$ là $b\alpha_2$.
\item \textbf{Cột $c$:} Ta được véc tơ $\alpha_3 = (1; 0; 0; 1; 0) \implies$ phần chứa $c$ là $c\alpha_3$.
\item \textbf{Cột $d$:} Ta được véc tơ $\alpha_4 = (-2; 0; 0; 0; 1) \implies$ phần chứa $d$ là $d\alpha_4$.
\end{itemize}
$\Rightarrow$ Tách thành $X = a\alpha_1 + b\alpha_2 + c\alpha_3 + d\alpha_4$ cực kỳ nhanh và không bao giờ bị nhầm dấu!
\end{quote}

Đặt $\alpha_1 = (2; 1; 0; 0; 0)$, $\alpha_2 = (-3; 0; 1; 0; 0)$, $\alpha_3 = (1; 0; 0; 1; 0)$, $\alpha_4 = (-2; 0; 0; 0; 1)$ và gọi $S = \{ \alpha_1, \alpha_2, \alpha_3, \alpha_4 \}$.

$\Rightarrow W = \{ X = a\alpha_1 + b\alpha_2 + c\alpha_3 + d\alpha_4 \mid a, b, c, d \in \R; \alpha_i \in S \} = \langle S \rangle$.

Mà $\langle S \rangle \le V$ nên $W = \langle S \rangle \le V$ (đpcm).

Ta có $W = \langle S \rangle$ nên $S = \{ \alpha_1, \alpha_2, \alpha_3, \alpha_4 \}$ là một hệ sinh của $W$.

Tiếp theo, ta tìm cơ sở và số chiều cho $W$ như sau:

Lập ma trận
\[
A_S = \begin{pmatrix} \alpha_1 \\ \alpha_2 \\ \alpha_3 \\ \alpha_4 \end{pmatrix} = 
\begin{pmatrix} 
2 & 1 & 0 & 0 & 0 \\ 
-3 & 0 & 1 & 0 & 0 \\ 
1 & 0 & 0 & 1 & 0 \\ 
-2 & 0 & 0 & 0 & 1 
\end{pmatrix}
\xrightarrow{(1) \leftrightarrow (3)}
\begin{pmatrix} 
1 & 0 & 0 & 1 & 0 \\ 
-3 & 0 & 1 & 0 & 0 \\ 
2 & 1 & 0 & 0 & 0 \\ 
-2 & 0 & 0 & 0 & 1 
\end{pmatrix}
\]
\[
\xrightarrow{\substack{(2) \to (2) + 3(1) \\ (3) \to (3) - 2(1) \\ (4) \to (4) + 2(1)}}
\begin{pmatrix} 
1 & 0 & 0 & 1 & 0 \\ 
0 & 0 & 1 & 3 & 0 \\ 
0 & 1 & 0 & -2 & 0 \\ 
0 & 0 & 0 & 2 & 1 
\end{pmatrix}
\xrightarrow{(2) \leftrightarrow (3)}
\begin{pmatrix} 
1 & 0 & 0 & 1 & 0 \\ 
0 & 1 & 0 & -2 & 0 \\ 
0 & 0 & 1 & 3 & 0 \\ 
0 & 0 & 0 & 2 & 1 
\end{pmatrix}
\]
Do ma trận đáp số có 4 dòng khác zero nên $r(A_S) = 4 = $ số véc tơ trong $S$, nên $S$ vừa là hệ sinh, vừa ĐLTT nên $S$ cũng là 1 cơ sở của $W$, và số chiều là $\dim_{\R} W = 4$.

\vspace{0.3cm}

\subsection*{BÀI 2 (Mẫu Câu 2 trên $V = \R^3$):}
Trên $V=\R^3$, cho tập hợp
\[
a=\{u_1=(1; 1; 1); u_2=(1; 1; 0); u_3=(1; 0; 0)\}.
\]
\begin{enumerate}[label=\alph*/,noitemsep]
\item Chứng minh rằng $a$ là một cơ sở của $\R^3$.
\item Trực chuẩn hóa cơ sở $a$ thành cơ sở trực chuẩn $b$.
\item Tìm tọa độ của một véc tơ $v=(2; 3; 4) \in \R^3$ theo cơ sở $a$.
\item Tìm tọa độ của véc tơ $v=(2; 3; 4)$ theo cơ sở $b$.
\end{enumerate}

\underline{\textbf{Giải:}}

\textbf{a/} Do $\R^3$ là không gian tuyến tính 3 chiều trên trường số thực $\R$, nên mỗi cơ sở của $\R^3$ là một hệ sinh, có 3 véc tơ, ĐLTT. Mà tập hợp $a$ đã có 3 véc tơ nên $a$ là 1 hệ sinh của $\R^3$. Để chứng tỏ $a$ là một cơ sở của $\R^3$ thì ta chỉ cần chứng tỏ $a$ là ĐLTT.

Ta lập ma trận
\[
A_a = \begin{pmatrix} u_1 \\ u_2 \\ u_3 \end{pmatrix} = \begin{pmatrix} 1 & 1 & 1 \\ 1 & 1 & 0 \\ 1 & 0 & 0 \end{pmatrix}
\]
Ta có $\det(A_a) = 1(0 - 0) - 1(0 - 0) + 1(0 - 1) = -1 \neq 0$.

$\Rightarrow a$ là một hệ sinh ĐLTT, nên $a$ là một cơ sở của $\R^3$.

\begin{quote}
\small
\noindent\fbox{\textbf{NOTE / GIẢI THÍCH CÂU A (Cách 2 -- Biến đổi sơ cấp dòng)}}\par\vspace{3pt}
Ta có thể đưa ma trận về bậc thang:
\[
A_a = \begin{pmatrix} 1 & 1 & 1 \\ 1 & 1 & 0 \\ 1 & 0 & 0 \end{pmatrix} \xrightarrow{\substack{(2) \to (2) - (1) \\ (3) \to (3) - (1)}} \begin{pmatrix} 1 & 1 & 1 \\ 0 & 0 & -1 \\ 0 & -1 & -1 \end{pmatrix} \xrightarrow{(2) \leftrightarrow (3)} \begin{pmatrix} 1 & 1 & 1 \\ 0 & -1 & -1 \\ 0 & 0 & -1 \end{pmatrix}
\]
Ma trận bậc thang có 3 dòng khác zero $\implies r(A_a) = 3 =$ số véc tơ trong $a \implies a$ ĐLTT $\implies a$ là một cơ sở của $\R^3$.
\end{quote}

\vspace{0.2cm}
\textbf{b/} Dùng pp Gram--Schmidt, ta làm như sau:

\begin{quote}
\small
\noindent\fbox{\textbf{NOTE / QUY TẮC TÍCH VÔ HƯỚNG TIÊU CHUẨN}}\par\vspace{3pt}
Đề bài không nói gì thêm nghĩa là dùng tích vô hướng chính tắc trong $\R^3$: $\ip{x}{y} = x_1y_1 + x_2y_2 + x_3y_3$ (lấy từng tọa độ tương ứng nhân với nhau rồi cộng lại). Độ dài bình phương $\norm{x}^2 = x_1^2 + x_2^2 + x_3^2$ (bình phương từng tọa độ rồi cộng lại).
\end{quote}

\textbf{Bước 1: Tìm $v_1$:}
\[
v_1 = u_1 = (1, 1, 1).
\]

\begin{quote}
\small
\noindent\fbox{\textbf{NOTE / GIẢI THÍCH BƯỚC 1}}\par\vspace{3pt}
Tính sẵn mẫu số chuẩn bị cho bước sau: $\norm{v_1}^2 = 1^2 + 1^2 + 1^2 = \mathbf{3}$.
\end{quote}

\vspace{0.2cm}
\textbf{Bước 2: Tìm $v_2$:}
\begin{align*}
v_2 &= u_2 - \frac{\ip{u_2}{v_1}}{\norm{v_1}^2}v_1 \\
    &= (1, 1, 0) - \frac{1(1) + 1(1) + 0(1)}{1^2 + 1^2 + 1^2}(1, 1, 1) \\
    &= (1, 1, 0) - \frac{2}{3}(1, 1, 1) = \left(\frac{1}{3}, \frac{1}{3}, -\frac{2}{3}\right).
\end{align*}

\begin{quote}
\small
\noindent\fbox{\textbf{NOTE / GIẢI THÍCH BƯỚC 2 (Cách ráp số)}}\par\vspace{3pt}
\begin{itemize}[noitemsep,topsep=2pt]
\item \textbf{Tử số $\ip{u_2}{v_1}$}: Ghép $u_2 = (\mathbf{1}; \mathbf{1}; \mathbf{0})$ với $v_1 = (\mathbf{1}; \mathbf{1}; \mathbf{1}) \implies 1(1) + 1(1) + 0(1) = 1 + 1 + 0 = \mathbf{2}$.
\item \textbf{Mẫu số}: $\norm{v_1}^2 = 1^2 + 1^2 + 1^2 = \mathbf{3} \implies$ hệ số là $\dfrac{\mathbf{2}}{\mathbf{3}}$.
\item \textbf{Phép trừ}: $(1; 1; 0) - \left(\dfrac{2}{3}; \dfrac{2}{3}; \dfrac{2}{3}\right) = \left(1-\dfrac{2}{3}; 1-\dfrac{2}{3}; 0-\dfrac{2}{3}\right) = \left(\dfrac{1}{3}; \dfrac{1}{3}; -\dfrac{2}{3}\right)$.
\item \textbf{Mẫu số cho bước sau}: $\norm{v_2}^2 = \left(\dfrac{1}{3}\right)^2 + \left(\dfrac{1}{3}\right)^2 + \left(-\dfrac{2}{3}\right)^2 = \dfrac{1}{9} + \dfrac{1}{9} + \dfrac{4}{9} = \dfrac{6}{9} = \mathbf{\dfrac{2}{3}}$.
\end{itemize}
\end{quote}

\vspace{0.2cm}
\textbf{Bước 3: Tìm $v_3$:}
\begin{align*}
v_3 &= u_3 - \frac{\ip{u_3}{v_2}}{\norm{v_2}^2}v_2 - \frac{\ip{u_3}{v_1}}{\norm{v_1}^2}v_1 \\
    &= (1, 0, 0) - \frac{1(1/3) + 0(1/3) + 0(-2/3)}{(1/3)^2 + (1/3)^2 + (-2/3)^2}\left(\frac{1}{3}, \frac{1}{3}, -\frac{2}{3}\right) - \frac{1(1) + 0(1) + 0(1)}{3}(1, 1, 1) \\
    &= (1, 0, 0) - \frac{1/3}{6/9}\left(\frac{1}{3}, \frac{1}{3}, -\frac{2}{3}\right) - \frac{1}{3}(1, 1, 1) \\
    &= (1, 0, 0) - \frac{1}{2}\left(\frac{1}{3}, \frac{1}{3}, -\frac{2}{3}\right) - \left(\frac{1}{3}, \frac{1}{3}, \frac{1}{3}\right) = \left(\frac{1}{2}, -\frac{1}{2}, 0\right).
\end{align*}

\begin{quote}
\small
\noindent\fbox{\textbf{NOTE / GIẢI THÍCH BƯỚC 3 (Cách ráp số)}}\par\vspace{3pt}
\begin{itemize}[noitemsep,topsep=2pt]
\item \textbf{Cụm $v_2$}: Ghép $u_3 = (\mathbf{1}; \mathbf{0}; \mathbf{0})$ với $v_2 = \left(\mathbf{\dfrac{1}{3}}; \mathbf{\dfrac{1}{3}}; \mathbf{-\dfrac{2}{3}}\right) \implies \ip{u_3}{v_2} = 1\left(\dfrac{1}{3}\right) + 0\left(\dfrac{1}{3}\right) + 0\left(-\dfrac{2}{3}\right) = \dfrac{1}{3}$. Mẫu $\norm{v_2}^2 = \dfrac{6}{9} \implies$ hệ số $\dfrac{1/3}{6/9} = \mathbf{\dfrac{1}{2}}$.
\item \textbf{Cụm $v_1$}: Ghép $u_3 = (\mathbf{1}; \mathbf{0}; \mathbf{0})$ với $v_1 = (\mathbf{1}; \mathbf{1}; \mathbf{1}) \implies \ip{u_3}{v_1} = 1(1) + 0(1) + 0(1) = 1$. Mẫu $\norm{v_1}^2 = 3 \implies$ hệ số $\mathbf{\dfrac{1}{3}}$.
\item \textbf{Phép trừ}: $(1; 0; 0) - \left(\dfrac{1}{6}; \dfrac{1}{6}; -\dfrac{1}{3}\right) - \left(\dfrac{1}{3}; \dfrac{1}{3}; \dfrac{1}{3}\right) = \left(\dfrac{1}{2}; -\dfrac{1}{2}; 0\right)$.
\end{itemize}
\end{quote}

\vspace{0.2cm}
\textbf{Bước 4: Chuẩn hóa thành cơ sở trực chuẩn $b$:}
\[
q_1 = \frac{v_1}{\norm{v_1}} = \frac{1}{\sqrt{3}}(1, 1, 1) = \left(\frac{1}{\sqrt{3}}, \frac{1}{\sqrt{3}}, \frac{1}{\sqrt{3}}\right),
\]
\[
q_2 = \frac{v_2}{\norm{v_2}} = \frac{1}{\sqrt{6/9}}\left(\frac{1}{3}, \frac{1}{3}, -\frac{2}{3}\right) = \frac{3}{\sqrt{6}}\left(\frac{1}{3}, \frac{1}{3}, -\frac{2}{3}\right) = \left(\frac{1}{\sqrt{6}}, \frac{1}{\sqrt{6}}, -\frac{2}{\sqrt{6}}\right),
\]
\[
q_3 = \frac{v_3}{\norm{v_3}} = \frac{1}{\sqrt{2/4}}\left(\frac{1}{2}, -\frac{1}{2}, 0\right) = \frac{2}{\sqrt{2}}\left(\frac{1}{2}, -\frac{1}{2}, 0\right) = \left(\frac{1}{\sqrt{2}}, -\frac{1}{\sqrt{2}}, 0\right).
\]
Suy ra $b = \{q_1, q_2, q_3\}$ là một cơ sở trực chuẩn cần tìm.

\begin{quote}
\small
\noindent\fbox{\textbf{NOTE / GIẢI THÍCH BƯỚC 4 (Liệt kê $v_i$ và độ dài $\norm{v_i}$ để ráp số):}}\par\vspace{3pt}
\begin{itemize}[noitemsep,topsep=2pt]
\item Với $q_1$: $v_1 = (1, 1, 1)$ và độ dài $\norm{v_1} = \sqrt{1^2 + 1^2 + 1^2} = \mathbf{\sqrt{3}}$.
\item Với $q_2$: $v_2 = \left(\dfrac{1}{3}; \dfrac{1}{3}; -\dfrac{2}{3}\right)$ và độ dài $\norm{v_2} = \sqrt{\left(\dfrac{1}{3}\right)^2 + \left(\dfrac{1}{3}\right)^2 + \left(-\dfrac{2}{3}\right)^2} = \mathbf{\sqrt{\dfrac{6}{9}}} = \dfrac{\sqrt{6}}{3}$.
\item Với $q_3$: $v_3 = \left(\dfrac{1}{2}; -\dfrac{1}{2}; 0\right)$ và độ dài $\norm{v_3} = \sqrt{\left(\dfrac{1}{2}\right)^2 + \left(-\dfrac{1}{2}\right)^2 + 0^2} = \mathbf{\sqrt{\dfrac{2}{4}}} = \dfrac{\sqrt{2}}{2}$.
\end{itemize}
\end{quote}

\vspace{0.2cm}
\textbf{c/} Tìm tọa độ của $v=(2; 3; 4)$ theo cơ sở $a$:

Ta cần tìm $c_1, c_2, c_3 \in \R$ sao cho: $v = c_1u_1 + c_2u_2 + c_3u_3$
\begin{align*}
(2, 3, 4) &= c_1(1, 1, 1) + c_2(1, 1, 0) + c_3(1, 0, 0) \\
&\Leftrightarrow \begin{cases} 
c_1 + c_2 + c_3 = 2 \\ 
c_1 + c_2 = 3 \\ 
c_1 = 4 
\end{cases} \Leftrightarrow \begin{cases} 
c_1 = 4 \\ 
c_2 = -1 \\ 
c_3 = -1 
\end{cases}
\end{align*}
Như vậy, ta có: $[v]_a = \begin{pmatrix} c_1 \\ c_2 \\ c_3 \end{pmatrix} = \begin{pmatrix} 4 \\ -1 \\ -1 \end{pmatrix}$.

\begin{quote}
\small
\noindent\fbox{\textbf{NOTE / GIẢI THÍCH CÂU C (Cách lập và giải hệ tìm tọa độ):}}\par\vspace{3pt}
\begin{itemize}[noitemsep,topsep=2pt]
\item \textbf{Bản chất}: Tọa độ của $v$ theo cơ sở $a$ là bộ 3 số $(c_1, c_2, c_3)$ sao cho biểu diễn tuyến tính $v = c_1u_1 + c_2u_2 + c_3u_3$.
\item \textbf{Cách lập hệ}: Đồng nhất từng vị trí tọa độ của hai vế:
  \begin{itemize}[noitemsep]
  \item Vị trí 1: $c_1(1) + c_2(1) + c_3(1) = 2 \implies c_1 + c_2 + c_3 = 2$.
  \item Vị trí 2: $c_1(1) + c_2(1) + c_3(0) = 3 \implies c_1 + c_2 = 3$.
  \item Vị trí 3: $c_1(1) + c_2(0) + c_3(0) = 4 \implies c_1 = 4$.
  \end{itemize}
\item \textbf{Mẹo giải nhanh}: Hệ dạng bậc thang tam giác nên ta thế ngược từ dưới lên:
  \begin{itemize}[noitemsep]
  \item Phương trình cuối cho ngay $c_1 = 4$.
  \item Thế $c_1=4$ lên phương trình giữa: $4 + c_2 = 3 \implies c_2 = -1$.
  \item Thế $c_1=4, c_2=-1$ lên phương trình đầu: $4 - 1 + c_3 = 2 \implies c_3 = -1$.
  \end{itemize}
\item \textbf{Lưu ý cách viết}: Tọa độ $[v]_a$ luôn viết dưới dạng \textbf{ma trận cột}: $\begin{pmatrix} c_1 \\ c_2 \\ c_3 \end{pmatrix}$.
\end{itemize}
\end{quote}

\vspace{0.2cm}
\textbf{d/} Tìm tọa độ của $v=(2; 3; 4)$ theo cơ sở trực chuẩn $b$:

Do $b$ là cơ sở trực chuẩn của $\R^3$, nên ta có:
\[
d_1 = \ip{v}{q_1} = 2\left(\frac{1}{\sqrt{3}}\right) + 3\left(\frac{1}{\sqrt{3}}\right) + 4\left(\frac{1}{\sqrt{3}}\right) = \frac{9}{\sqrt{3}} = 3\sqrt{3},
\]
\[
d_2 = \ip{v}{q_2} = 2\left(\frac{1}{\sqrt{6}}\right) + 3\left(\frac{1}{\sqrt{6}}\right) + 4\left(-\frac{2}{\sqrt{6}}\right) = -\frac{3}{\sqrt{6}},
\]
\[
d_3 = \ip{v}{q_3} = 2\left(\frac{1}{\sqrt{2}}\right) + 3\left(-\frac{1}{\sqrt{2}}\right) + 4(0) = -\frac{1}{\sqrt{2}}.
\]
Suy ra $[v]_b = \begin{pmatrix} 3\sqrt{3} \\ -3/\sqrt{6} \\ -1/\sqrt{2} \end{pmatrix}$.

\begin{quote}
\small
\noindent\fbox{\textbf{NOTE / GIẢI THÍCH CÂU D (Công thức tọa độ trong cơ sở trực chuẩn):}}\par\vspace{3pt}
\begin{itemize}[noitemsep,topsep=2pt]
\item \textbf{Tại sao không cần giải hệ phương trình như câu c}: Do $b=\{q_1, q_2, q_3\}$ là cơ sở \textbf{TRỰC CHUẨN}, ta tính trực tiếp từng tọa độ bằng tích vô hướng: $d_i = \ip{v}{q_i}$ với $i=1, 2, 3$.
\item \textbf{Cách ráp số từng tọa độ}:
  \begin{itemize}[noitemsep]
  \item $d_1 = \ip{v}{q_1}$: Lấy $v=(2, 3, 4)$ nhân từng vị trí với $q_1 = \left(\dfrac{1}{\sqrt{3}}, \dfrac{1}{\sqrt{3}}, \dfrac{1}{\sqrt{3}}\right) \implies 2\left(\dfrac{1}{\sqrt{3}}\right) + 3\left(\dfrac{1}{\sqrt{3}}\right) + 4\left(\dfrac{1}{\sqrt{3}}\right) = \dfrac{9}{\sqrt{3}} = 3\sqrt{3}$.
  \item $d_2 = \ip{v}{q_2}$: Lấy $v=(2, 3, 4)$ nhân từng vị trí với $q_2 = \left(\dfrac{1}{\sqrt{6}}, \dfrac{1}{\sqrt{6}}, -\dfrac{2}{\sqrt{6}}\right) \implies 2\left(\dfrac{1}{\sqrt{6}}\right) + 3\left(\dfrac{1}{\sqrt{6}}\right) + 4\left(-\dfrac{2}{\sqrt{6}}\right) = \dfrac{2+3-8}{\sqrt{6}} = -\dfrac{3}{\sqrt{6}}$.
  \item $d_3 = \ip{v}{q_3}$: Lấy $v=(2, 3, 4)$ nhân từng vị trí với $q_3 = \left(\dfrac{1}{\sqrt{2}}, -\dfrac{1}{\sqrt{2}}, 0\right) \implies 2\left(\dfrac{1}{\sqrt{2}}\right) + 3\left(-\dfrac{1}{\sqrt{2}}\right) + 4(0) = \dfrac{2-3+0}{\sqrt{2}} = -\dfrac{1}{\sqrt{2}}$.
  \end{itemize}
\item \textbf{Kết luận}: Tọa độ $[v]_b$ viết theo cột: $[v]_b = \begin{pmatrix} d_1 \\ d_2 \\ d_3 \end{pmatrix} = \begin{pmatrix} 3\sqrt{3} \\ -3/\sqrt{6} \\ -1/\sqrt{2} \end{pmatrix}$.
\end{itemize}
\end{quote}


\vspace{0.3cm}

\subsection*{BÀI 3A (Mẫu Câu 3 trên $V = \R^4$ với tích vô hướng tiêu chuẩn):}
Trên $V=\R^4$, với tích vô hướng $\ip{x}{y} = x_1y_1+x_2y_2+x_3y_3+x_4y_4$ (tích vô hướng tiêu chuẩn), cho hệ
\[
S=\{u_1=(1, -1, 1, -1); u_2=(1, 1, 1, 1); u_3=(1, 0, 0, 0)\}.
\]
Dùng pp Gram--Schmidt để trực chuẩn hóa $S$. 
Cho véc tơ $\lambda=(2, 1, 4, 1)$. Tìm tọa độ của véc tơ $\lambda$ theo cơ sở vừa trực chuẩn hóa được.

\underline{\textbf{Giải:}}

Dùng pp Gram--Schmidt, ta làm như sau:
\[
v_1 = u_1 = (1, -1, 1, -1),
\]
\begin{align*}
v_2 &= u_2 - \frac{\ip{u_2}{v_1}}{\norm{v_1}^2}v_1 \\
    &= (1, 1, 1, 1) - \frac{1(1) + 1(-1) + 1(1) + 1(-1)}{4}(1, -1, 1, -1) \\
    &= (1, 1, 1, 1) - 0(1, -1, 1, -1) = (1, 1, 1, 1),
\end{align*}
\begin{align*}
v_3 &= u_3 - \frac{\ip{u_3}{v_2}}{\norm{v_2}^2}v_2 - \frac{\ip{u_3}{v_1}}{\norm{v_1}^2}v_1 \\
    &= (1, 0, 0, 0) - \frac{1(1)}{4}(1, 1, 1, 1) - \frac{1(1)}{4}(1, -1, 1, -1) \\
    &= (1, 0, 0, 0) - \left(\frac{1}{4}, \frac{1}{4}, \frac{1}{4}, \frac{1}{4}\right) - \left(\frac{1}{4}, -\frac{1}{4}, \frac{1}{4}, -\frac{1}{4}\right) \\
    &= \left(\frac{1}{2}, 0, -\frac{1}{2}, 0\right).
\end{align*}

\begin{quote}
\small
\noindent\fbox{\textbf{NOTE / GIẢI THÍCH (Trực giao hóa Gram--Schmidt):}}\par\vspace{3pt}
\begin{itemize}[noitemsep,topsep=2pt]
\item \textbf{Mẫu số $\norm{v_1}^2$}: $1^2 + (-1)^2 + 1^2 + (-1)^2 = \mathbf{4}$.
\item \textbf{Tính $v_2$}: Tử số $\ip{u_2}{v_1} = 1(1) + 1(-1) + 1(1) + 1(-1) = 0 \implies$ hệ số bằng $0$. Do đó $v_2 = u_2$ (vì $u_2$ đã vuông góc sẵn với $v_1$). Mẫu số $\norm{v_2}^2 = 1^2 + 1^2 + 1^2 + 1^2 = \mathbf{4}$.
\item \textbf{Tính $v_3$}: Cụm $v_2$ có tử số $\ip{u_3}{v_2} = 1(1) + 0 + 0 + 0 = 1$. Cụm $v_1$ có tử số $\ip{u_3}{v_1} = 1(1) + 0 + 0 + 0 = 1$. Cả hai cụm đều có hệ số $\dfrac{1}{4}$. Phép trừ: $(1, 0, 0, 0) - (1/4, 1/4, 1/4, 1/4) - (1/4, -1/4, 1/4, -1/4) = (1/2, 0, -1/2, 0)$.
\item \textbf{Mẫu số $\norm{v_3}^2$}: $(1/2)^2 + 0^2 + (-1/2)^2 + 0^2 = 1/4 + 1/4 = \mathbf{1/2}$.
\end{itemize}
\end{quote}

\vspace{0.2cm}
Ta trực chuẩn hóa tiếp theo:
\[
q_1 = \frac{v_1}{\norm{v_1}} = \frac{1}{\sqrt{4}}(1, -1, 1, -1) = \left(\frac{1}{2}, -\frac{1}{2}, \frac{1}{2}, -\frac{1}{2}\right),
\]
\[
q_2 = \frac{v_2}{\norm{v_2}} = \frac{1}{\sqrt{4}}(1, 1, 1, 1) = \left(\frac{1}{2}, \frac{1}{2}, \frac{1}{2}, \frac{1}{2}\right),
\]
\[
q_3 = \frac{v_3}{\norm{v_3}} = \frac{1}{\sqrt{1/4 + 1/4}}\left(\frac{1}{2}, 0, -\frac{1}{2}, 0\right) = \sqrt{2}\left(\frac{1}{2}, 0, -\frac{1}{2}, 0\right) = \left(\frac{1}{\sqrt{2}}, 0, -\frac{1}{\sqrt{2}}, 0\right).
\]
Tập hợp $b = \{q_1, q_2, q_3\}$ là cơ sở trực chuẩn vừa tìm được.

\begin{quote}
\small
\noindent\fbox{\textbf{NOTE / GIẢI THÍCH (Chuẩn hóa cơ sở trực chuẩn):}}\par\vspace{3pt}
\begin{itemize}[noitemsep,topsep=2pt]
\item $v_1 = (1, -1, 1, -1)$ có độ dài $\norm{v_1} = \sqrt{4} = 2 \implies q_1 = \dfrac{1}{2}v_1$.
\item $v_2 = (1, 1, 1, 1)$ có độ dài $\norm{v_2} = \sqrt{4} = 2 \implies q_2 = \dfrac{1}{2}v_2$.
\item $v_3 = (1/2, 0, -1/2, 0)$ có độ dài $\norm{v_3} = \sqrt{1/2} = \dfrac{1}{\sqrt{2}} \implies q_3 = \dfrac{v_3}{1/\sqrt{2}} = \sqrt{2}v_3 = \left(\dfrac{1}{\sqrt{2}}, 0, -\dfrac{1}{\sqrt{2}}, 0\right)$.
\end{itemize}
\end{quote}

\vspace{0.2cm}
Vì $b$ là cơ sở trực chuẩn, tọa độ của $\lambda = (2, 1, 4, 1)$ theo $b$ được tính như sau:
\[
c_1 = \ip{\lambda}{q_1} = 2\left(\frac{1}{2}\right) + 1\left(-\frac{1}{2}\right) + 4\left(\frac{1}{2}\right) + 1\left(-\frac{1}{2}\right) = 1 - \frac{1}{2} + 2 - \frac{1}{2} = 2,
\]
\[
c_2 = \ip{\lambda}{q_2} = 2\left(\frac{1}{2}\right) + 1\left(\frac{1}{2}\right) + 4\left(\frac{1}{2}\right) + 1\left(\frac{1}{2}\right) = 1 + \frac{1}{2} + 2 + \frac{1}{2} = 4,
\]
\[
c_3 = \ip{\lambda}{q_3} = 2\left(\frac{1}{\sqrt{2}}\right) + 1(0) + 4\left(-\frac{1}{\sqrt{2}}\right) + 1(0) = -\frac{2}{\sqrt{2}} = -\sqrt{2}.
\]
Suy ra $[\lambda]_b = \begin{pmatrix} 2 \\ 4 \\ -\sqrt{2} \end{pmatrix}$.

\begin{quote}
\small
\noindent\fbox{\textbf{NOTE / GIẢI THÍCH (Tính tọa độ qua tích vô hướng):}}\par\vspace{3pt}
Do $b$ là cơ sở \textbf{TRỰC CHUẨN}, tọa độ thứ $i$ được tính trực tiếp bằng tích vô hướng: $c_i = \ip{\lambda}{q_i}$ (không cần lập hệ phương trình). Kết luận viết theo ma trận cột: $[\lambda]_b = \begin{pmatrix} c_1 \\ c_2 \\ c_3 \end{pmatrix}$.
\end{quote}

\vspace{0.3cm}

\subsection*{BÀI 3B (Mẫu Câu 3 trên $V = \R^4$ với tích vô hướng có trọng số):}
Trên $V=\R^4$, với tích vô hướng $\ip{x}{y} = x_1y_1 + 2x_2y_2 + x_3y_3 + 2x_4y_4$, cho hệ
\[
S=\{u_1=(1; 1; 0; 0); u_2=(1; 0; 1; 1); u_3=(0; 1; 1; 0)\}.
\]
Dùng pp Gram--Schmidt để trực chuẩn hóa $S$.
Cho $\lambda = (3; 1; 2; -1)$. Tìm tọa độ của véc tơ $\lambda$ theo cơ sở vừa trực chuẩn hóa được.

\underline{\textbf{Giải:}}

Ta tính theo tích vô hướng đã cho:
\[
v_1 = u_1 = (1; 1; 0; 0),
\]
\[
\norm{v_1}^2 = \ip{v_1}{v_1} = 1^2 + 2(1^2) + 0^2 + 2(0^2) = 1 + 2 = 3.
\]
\begin{align*}
v_2 &= u_2 - \frac{\ip{u_2}{v_1}}{\norm{v_1}^2}v_1 \\
    &= (1; 0; 1; 1) - \frac{1(1) + 2(0)(1) + 1(0) + 2(1)(0)}{3}(1; 1; 0; 0) \\
    &= (1; 0; 1; 1) - \frac{1}{3}(1; 1; 0; 0) = \left(\frac{2}{3}; -\frac{1}{3}; 1; 1\right).
\end{align*}
Tính $\norm{v_2}^2$:
\[
\norm{v_2}^2 = \left(\frac{2}{3}\right)^2 + 2\left(-\frac{1}{3}\right)^2 + 1^2 + 2(1^2) = \frac{4}{9} + \frac{2}{9} + 1 + 2 = \frac{6}{9} + 3 = \frac{11}{3}.
\]
Tiếp theo, ta tính $v_3$:
\[
\ip{u_3}{v_1} = 0(1) + 2(1)(1) + 1(0) + 2(0)(0) = 2,
\]
\[
\ip{u_3}{v_2} = 0\left(\frac{2}{3}\right) + 2(1)\left(-\frac{1}{3}\right) + 1(1) + 2(0)(1) = -\frac{2}{3} + 1 = \frac{1}{3}.
\]
\begin{align*}
v_3 &= u_3 - \frac{\ip{u_3}{v_2}}{\norm{v_2}^2}v_2 - \frac{\ip{u_3}{v_1}}{\norm{v_1}^2}v_1 \\
    &= (0; 1; 1; 0) - \frac{1/3}{11/3}\left(\frac{2}{3}; -\frac{1}{3}; 1; 1\right) - \frac{2}{3}(1; 1; 0; 0) \\
    &= (0; 1; 1; 0) - \frac{1}{11}\left(\frac{2}{3}; -\frac{1}{3}; 1; 1\right) - \left(\frac{2}{3}; \frac{2}{3}; 0; 0\right) \\
    &= \left(-\frac{2}{33} - \frac{2}{3}; 1 + \frac{1}{33} - \frac{2}{3}; 1 - \frac{1}{11}; -\frac{1}{11}\right) = \left(-\frac{8}{11}; \frac{4}{11}; \frac{10}{11}; -\frac{1}{11}\right).
\end{align*}
Tính độ dài $\norm{v_3}^2$:
\[
\norm{v_3}^2 = \left(-\frac{8}{11}\right)^2 + 2\left(\frac{4}{11}\right)^2 + \left(\frac{10}{11}\right)^2 + 2\left(-\frac{1}{11}\right)^2 = \frac{64 + 32 + 100 + 2}{121} = \frac{198}{121} = \frac{18}{11}.
\]

\begin{quote}
\small
\noindent\fbox{\textbf{NOTE / GIẢI THÍCH (Tích vô hướng CÓ TRỌNG SỐ):}}\par\vspace{3pt}
\begin{itemize}[noitemsep,topsep=2pt]
\item \textbf{Cực kỳ quan trọng}: Đề bài cho $\ip{x}{y} = x_1y_1 + \mathbf{2}x_2y_2 + x_3y_3 + \mathbf{2}x_4y_4$. Mọi phép tính tích vô hướng và độ dài $\norm{x}^2$ đều \textbf{phải nhân thêm hệ số 2} ở vị trí tọa độ thứ hai và thứ tư:
  $$\norm{x}^2 = x_1^2 + \mathbf{2}x_2^2 + x_3^2 + \mathbf{2}x_4^2$$
\item $\norm{v_1}^2 = 1^2 + \mathbf{2}(1^2) + 0 + 0 = \mathbf{3}$.
\item Tử số $\ip{u_2}{v_1} = 1(1) + \mathbf{2}(0)(1) + 1(0) + \mathbf{2}(1)(0) = \mathbf{1} \implies$ hệ số là $\dfrac{1}{3}$.
\item $\norm{v_2}^2 = (2/3)^2 + \mathbf{2}(-1/3)^2 + 1^2 + \mathbf{2}(1^2) = 4/9 + 2/9 + 3 = \mathbf{11/3}$.
\item Với $v_3$: Hệ số cụm $v_2$ là $\dfrac{1/3}{11/3} = \mathbf{\dfrac{1}{11}}$; hệ số cụm $v_1$ là $\mathbf{\dfrac{2}{3}}$.
\item $\norm{v_3}^2 = (-8/11)^2 + \mathbf{2}(4/11)^2 + (10/11)^2 + \mathbf{2}(-1/11)^2 = \dfrac{64+32+100+2}{121} = \mathbf{\dfrac{18}{11}}$.
\end{itemize}
\end{quote}

\vspace{0.2cm}
Trực chuẩn hóa:
\[
q_1 = \frac{v_1}{\norm{v_1}} = \frac{1}{\sqrt{3}}(1; 1; 0; 0),
\]
\[
q_2 = \frac{v_2}{\norm{v_2}} = \frac{1}{\sqrt{11/3}}\left(\frac{2}{3}; -\frac{1}{3}; 1; 1\right) = \frac{\sqrt{3}}{\sqrt{11}}\left(\frac{2}{3}; -\frac{1}{3}; 1; 1\right) = \frac{1}{\sqrt{33}}(2; -1; 3; 3),
\]
\[
q_3 = \frac{v_3}{\norm{v_3}} = \frac{1}{\sqrt{18/11}}\frac{1}{11}(-8; 4; 10; -1) = \frac{1}{3\sqrt{2}\sqrt{11}}(-8; 4; 10; -1) = \frac{1}{3\sqrt{22}}(-8; 4; 10; -1).
\]
Tập hợp $b = \{q_1, q_2, q_3\}$ là cơ sở trực chuẩn cần tìm.

\begin{quote}
\small
\noindent\fbox{\textbf{NOTE / GIẢI THÍCH (Cách rút gọn căn khi chuẩn hóa):}}\par\vspace{3pt}
\begin{itemize}[noitemsep,topsep=2pt]
\item Với $q_2$: $\norm{v_2} = \sqrt{11/3} \implies \dfrac{1}{\norm{v_2}} = \dfrac{\sqrt{3}}{\sqrt{11}}$. Khi nhân với $v_2 = \dfrac{1}{3}(2; -1; 3; 3)$, ta rút gọn $\dfrac{\sqrt{3}}{3\sqrt{11}} = \dfrac{1}{\sqrt{3}\sqrt{11}} = \mathbf{\dfrac{1}{\sqrt{33}}}$.
\item Với $q_3$: $\norm{v_3} = \sqrt{18/11} = \dfrac{3\sqrt{2}}{\sqrt{11}} \implies \dfrac{1}{\norm{v_3}} = \dfrac{\sqrt{11}}{3\sqrt{2}}$. Nhân với $v_3 = \dfrac{1}{11}(-8; 4; 10; -1)$, ta rút gọn $\dfrac{\sqrt{11}}{11 \cdot 3\sqrt{2}} = \mathbf{\dfrac{1}{3\sqrt{22}}}$.
\end{itemize}
\end{quote}

\vspace{0.2cm}
Tọa độ của $\lambda = (3; 1; 2; -1)$ theo cơ sở $b$:
\[
c_1 = \ip{\lambda}{q_1} = \frac{1}{\sqrt{3}}[3(1) + 2(1)(1) + 2(0) + 2(-1)(0)] = \frac{5}{\sqrt{3}},
\]
\[
c_2 = \ip{\lambda}{q_2} = \frac{1}{\sqrt{33}}[3(2) + 2(1)(-1) + 2(3) + 2(-1)(3)] = \frac{1}{\sqrt{33}}[6 - 2 + 6 - 6] = \frac{4}{\sqrt{33}},
\]
\begin{align*}
c_3 = \ip{\lambda}{q_3} &= \frac{1}{3\sqrt{22}}[3(-8) + 2(1)(4) + 2(10) + 2(-1)(-1)] \\
&= \frac{1}{3\sqrt{22}}[-24 + 8 + 20 + 2] = \frac{6}{3\sqrt{22}} = \frac{2}{\sqrt{22}}.
\end{align*}
Suy ra $[\lambda]_b = \begin{pmatrix} 5/\sqrt{3} \\ 4/\sqrt{33} \\ 2/\sqrt{22} \end{pmatrix}$.

\begin{quote}
\small
\noindent\fbox{\textbf{NOTE / GIẢI THÍCH (Tính tọa độ có trọng số):}}\par\vspace{3pt}
\textbf{Lỗi sai phổ biến}: Khi tính $c_i = \ip{\lambda}{q_i}$, phải áp dụng đúng công thức tích vô hướng có trọng số của đề bài (nhớ nhân 2 ở vị trí 2 và 4 của véc tơ $\lambda$ và $q_i$).
\end{quote}


\vspace{0.3cm}

\subsection*{BÀI 4 (Mẫu Câu 4 - Chéo hóa ma trận vuông cấp 2):}
Cho ma trận vuông thực, cấp $2$, $A = \begin{pmatrix} 4 & 2 \\ 3 & -1 \end{pmatrix}$. Hãy chéo hóa ma trận $A$, rồi tìm $A^m$, $\forall m\in\N$ và suy ra $A^{2026}$.

\underline{\textbf{Giải:}}

Ta có đa thức đặc trưng của $A$:
\[
p_A(\lambda) = \det(A - \lambda I_2) = \begin{vmatrix} 4 - \lambda & 2 \\ 3 & -1 - \lambda \end{vmatrix} = (4 - \lambda)(-1 - \lambda) - 6 = \lambda^2 - 3\lambda - 10
\]
Giải pt: $p_A(\lambda) = 0 \Leftrightarrow \lambda^2 - 3\lambda - 10 = 0 \Leftrightarrow \begin{bmatrix} \lambda = 5 \\ \lambda = -2 \end{bmatrix}$.

Do đó, $A$ có 2 trị riêng: $\lambda_1 = 5, \lambda_2 = -2$ và số mũ là $r_1=1, r_2=1$.

\begin{quote}
\small
\noindent\fbox{\textbf{NOTE / GIẢI THÍCH (Đa thức đặc trưng và Trị riêng):}}\par\vspace{3pt}
\begin{itemize}[noitemsep,topsep=2pt]
\item \textbf{Đa thức đặc trưng}: $p_A(\lambda) = \det(A - \lambda I_2) = (a_{11}-\lambda)(a_{22}-\lambda) - a_{12}a_{21}$.
\item Bấm máy tính phương trình bậc hai được 2 nghiệm $\lambda_1 = 5, \lambda_2 = -2$.
\item \textbf{Quy tắc Thầy giảng}: Ma trận vuông cấp 2 có 2 trị riêng thực phân biệt thì \textbf{chắc chắn chéo hóa được trên $\R$}, số mũ (bội đại số) của mỗi trị riêng đều bằng 1: $r_1 = 1, r_2 = 1$.
\end{itemize}
\end{quote}

\vspace{0.2cm}
Tiếp theo, ta tìm cơ sở $a_j$ cho các không gian riêng:

Với $\lambda_1 = 5$:
$E_5 = \{ X \in \R^2 \mid (A - 5I_2)X = O \}$

Ta có dạng ma trận hóa:
\[
\begin{bmatrix} -1 & 2 \mid 0 \\ 3 & -6 \mid 0 \end{bmatrix} \xrightarrow{(2) \to (2) + 3(1)} \begin{bmatrix} -1 & 2 \mid 0 \\ 0 & 0 \mid 0 \end{bmatrix}
\]
Hệ phương trình tương đương: $-x_1 + 2x_2 = 0 \Rightarrow x_1 = 2x_2$.

Đặt $x_2 = a \ (\forall a \in \R) \Rightarrow x_1 = 2a$.

$E_5 = \{ X = (2a, a) \mid a \in \R \} = \{ X = a(2, 1) \mid a \in \R \}$.

Đặt $u_1 = (2, 1)$, ta có $a_1 = \{u_1\}$ là cơ sở của $E_5$ và $\dim_{\R} E_5 = 1 = r_1$. \hfill (1)

Với $\lambda_2 = -2$:
$E_{-2} = \{ X \in \R^2 \mid (A + 2I_2)X = O \}$

Ta có dạng ma trận hóa:
\[
\begin{bmatrix} 6 & 2 \mid 0 \\ 3 & 1 \mid 0 \end{bmatrix} \xrightarrow{(2) \to (2) - \frac{1}{2}(1)} \begin{bmatrix} 6 & 2 \mid 0 \\ 0 & 0 \mid 0 \end{bmatrix}
\]
Hệ phương trình tương đương: $6x_1 + 2x_2 = 0 \Rightarrow x_2 = -3x_1$.

Đặt $x_1 = b \ (\forall b \in \R) \Rightarrow x_2 = -3b$.

$E_{-2} = \{ X = (b, -3b) \mid b \in \R \} = \{ X = b(1, -3) \mid b \in \R \}$.

Đặt $u_2 = (1, -3)$, ta có $a_2 = \{u_2\}$ là cơ sở của $E_{-2}$ và $\dim_{\R} E_{-2} = 1 = r_2$. \hfill (2)

Từ (1), (2) $\Rightarrow A$ có thể chéo hóa được trên $\R$.

\begin{quote}
\small
\noindent\fbox{\textbf{NOTE / GIẢI THÍCH (Tìm không gian riêng và véc tơ riêng):}}\par\vspace{3pt}
\begin{itemize}[noitemsep,topsep=2pt]
\item \textbf{Phương pháp}: Lập ma trận bổ sung $(A - \lambda I_2 \mid 0)$ rồi biến đổi sơ cấp dòng về bậc thang. Dòng thứ 2 sẽ luôn biến thành dòng toàn số $0$.
\item \textbf{Chọn ẩn tự do}:
  \begin{itemize}[noitemsep]
  \item Với $\lambda_1 = 5$: $x_1 = 2x_2$. Chọn $x_2 = 1 \implies x_1 = 2 \implies u_1 = (2, 1)$.
  \item Với $\lambda_2 = -2$: $x_2 = -3x_1$. Chọn $x_1 = 1 \implies x_2 = -3 \implies u_2 = (1, -3)$.
  \end{itemize}
\item \textbf{Lập luận ăn điểm}: $\dim_{\R} E_{\lambda_i} = r_i \implies A$ chéo hóa được trên $\R$.
\end{itemize}
\end{quote}

\vspace{0.2cm}
Đặt $a = a_1 \cup a_2 = \{u_1, u_2\}$ thì $a$ là một cơ sở của $\R^2$.

Gọi $\beta_0 = \{e_1 = (1, 0), e_2 = (0, 1)\}$ là cơ sở chính tắc của $\R^2$.

Đặt $T = T_{\beta_0 \to a} = (u_1 \quad u_2) = \begin{pmatrix} 2 & 1 \\ 1 & -3 \end{pmatrix}$.

Ta có: $D = T^{-1}AT = \begin{pmatrix} 5 & 0 \\ 0 & -2 \end{pmatrix}$. Đây là dạng chéo hóa cần tìm của $A$.

\begin{quote}
\small
\noindent\fbox{\textbf{NOTE / GIẢI THÍCH (Ghép ma trận $T$ và ma trận đường chéo $D$):}}\par\vspace{3pt}
\begin{itemize}[noitemsep,topsep=2pt]
\item $T$ được lập bằng cách ghép các véc tơ riêng làm cột: Cột 1 là $u_1$, cột 2 là $u_2$.
\item $D$ là ma trận đường chéo chứa các trị riêng \textbf{đúng theo thứ tự cột của $T$}: Cột 1 của $T$ là $u_1$ (ứng với $\lambda_1 = 5$) thì vị trí $(1, 1)$ của $D$ là số 5; cột 2 là $u_2$ (ứng với $\lambda_2 = -2$) thì vị trí $(2, 2)$ của $D$ là $-2$.
\end{itemize}
\end{quote}

\vspace{0.2cm}
Tìm $A^m$, $\forall m \in \N$:

Ta có $\det(T) = 2(-3) - 1(1) = -7 \neq 0$, suy ra:
\[
T^{-1} = \frac{1}{-7}\begin{pmatrix} -3 & -1 \\ -1 & 2 \end{pmatrix} = \frac{1}{7}\begin{pmatrix} 3 & 1 \\ 1 & -2 \end{pmatrix}.
\]
Lũy thừa:
\[
A^m = T D^m T^{-1} = \begin{pmatrix} 2 & 1 \\ 1 & -3 \end{pmatrix} \begin{pmatrix} 5^m & 0 \\ 0 & (-2)^m \end{pmatrix} \frac{1}{7}\begin{pmatrix} 3 & 1 \\ 1 & -2 \end{pmatrix}
\]
\[
= \frac{1}{7} \begin{pmatrix} 2\cdot 5^m & (-2)^m \\ 5^m & -3(-2)^m \end{pmatrix} \begin{pmatrix} 3 & 1 \\ 1 & -2 \end{pmatrix} 
= \frac{1}{7} \begin{pmatrix} 6\cdot 5^m + (-2)^m & 2\cdot 5^m - 2(-2)^m \\ 3\cdot 5^m - 3(-2)^m & 5^m + 6(-2)^m \end{pmatrix}, \quad \forall m \in \N.
\]
Suy ra $A^{2026}$ (với $m = 2026$, do $2026$ chẵn nên $(-2)^{2026} = 2^{2026}$):
\[
A^{2026} = \frac{1}{7} \begin{pmatrix} 6\cdot 5^{2026} + 2^{2026} & 2\cdot 5^{2026} - 2\cdot 2^{2026} \\ 3\cdot 5^{2026} - 3\cdot 2^{2026} & 5^{2026} + 6\cdot 2^{2026} \end{pmatrix}.
\]

\begin{quote}
\small
\noindent\fbox{\textbf{NOTE / GIẢI THÍCH (Tính lũy thừa $A^m$ và $A^{2026}$):}}\par\vspace{3pt}
\begin{itemize}[noitemsep,topsep=2pt]
\item \textbf{Nghịch đảo ma trận cấp 2 nhanh}:\\
Với $T = \begin{pmatrix} a & b \\ c & d \end{pmatrix} \implies T^{-1} = \dfrac{1}{\det(T)}\begin{pmatrix} d & -b \\ -c & a \end{pmatrix}$.
\item \textbf{Quy tắc nhân ma trận}: Giữ phân số $\dfrac{1}{7}$ ở ngoài. Nhân $T$ với $D^m$ trước (mỗi cột của $T$ nhân với phần tử đường chéo tương ứng của $D^m$), sau đó nhân với $T^{-1}$.
\item \textbf{Lưu ý số mũ chẵn}: Với $m = 2026$, vì 2026 là số chẵn nên $(-2)^{2026} = 2^{2026}$.
\end{itemize}
\end{quote}

\vspace{0.3cm}


\vspace{0.5cm}
\begin{center}
\vspace{0.3cm}
---------- HẾT ----------
\end{center}

\end{document}
